\chapter{Rularea proiectului}
\label{ch:rulare}

\paragraph{Cerințe.} Docker (cu Compose v2). Nu este nevoie de Rust instalat local: toate comenzile
\texttt{cargo} rulează într-un container (\texttt{scripts/cargo.sh}).

\begin{code}{Comenzi principale.}
make up            # docker compose up -d --build (11 servicii)
make consistency   # testul end-to-end de coerență (18 combinații)
make bench         # toate experimentele -> results/*.csv
make docs          # documentația PDF (docs/build/main.pdf)
make slides        # prezentarea PDF (presentation/build/slides.pdf)
make check         # fmt + clippy -D warnings + teste
make down          # oprește și șterge volumele
\end{code}

\paragraph{Adrese.}
\begin{itemize}
  \item interfața web: \url{http://localhost:8080}
  \item Swagger UI: \url{http://localhost:8080/swagger-ui/}
  \item Grafana: \url{http://localhost:3000}, Prometheus: \url{http://localhost:9090}
\end{itemize}

\paragraph{Schimbarea modului din linia de comandă.}
\begin{code}{}
curl -X POST localhost:8080/node/app1/admin/mode -H 'content-type: application/json' \
     -d '{"backend":"tiered","strategy":"write-through","capacity":50000,"policy":"clock"}'
\end{code}

\paragraph{O rulare de încărcare ad-hoc.}
\begin{code}{}
docker compose --profile bench run --rm loadgen http --concurrency 64 --zipf 0.99 \
     --read-ratio 0.95 --duration 20 --mode '{"backend":"cluster"}'
\end{code}
